\documentclass[twocolumn]{aastex631}
\begin{document}
\title{The reionization ionizing-photon-budget shortfall for star-forming galaxies is not robust to systematics at z\~{}5 (lowxi systematic corner)}
\author{NebulaMind Lab (autonomous pipeline)}
\affiliation{NebulaMind Lab, \url{https://lab.nebulamind.net}}
\begin{abstract}
Reionization ionizing-photon-budget reconciliation at z\~{}5: star-forming galaxies require f\_esc=0.049 (+0.050/-0.025) to close the budget (Madau-Dickinson SFRD, log xi\_ion=25.5+/-0.15, clumping C=2-5, JWST-SFRD tail), versus indirect-proxy-inferred f\_esc=0.062 (+0.108/-0.039) from LzLCS O32/beta calibrations. Median delta(required-inferred)=-0.011 dex-frac (16-84\%: -0.118 to +0.045); 42\% of the systematic MC shows a shortfall. Conclusion: the budget CLOSES within the systematic, and the sign holds under both O32 and beta calibrations. The result is bounded by the xi\_ion x clumping x proxy-calibration systematic, not by statistics. Generated autonomously from public data (jwst) via the NebulaMind Lab runner: a bounded, reproducible, descriptive study.
\end{abstract}
\section{Introduction}
Introduction: Recent studies have highlighted a potential crisis in reconciling the ionizing photon budget during reionization, with some suggesting that star-forming galaxies may not produce enough photons to account for the observed reionization [Muñoz2024]. This discrepancy has sparked interest in revisiting the assumptions and calibrations used in these calculations. To address this issue, we draw on previous work by Madau \& Dickinson (2014) and others, who have developed analytic models for cosmic star formation rate density (SFRD) and ionizing photon production.
\section{Data and method}
Data and method: Our analysis relies on published literature values for key parameters, including the Madau \& Dickinson (2014) SFRD fitting function, ionization fraction ξion, and escape fraction f\_esc proxy calibrations from LzLCS [Chisholm+22, Flury+22; Simmonds+24]. We do not use any new observational or catalog data in this study. Instead, we focus on systematically reconciling the reionization photon budget using these established values.
\section{Result}
Result: Our calculations show that star-forming galaxies can close the ionizing photon budget at z\~{}5 if they have an escape fraction f\_esc of 0.049 (+0.050/-0.025). This value is consistent with indirect-proxy-inferred estimates from LzLCS O32/beta calibrations, which yield f\_esc = 0.062 (+0.108/-0.039). The median difference between the required and inferred values is -0.011 dex-frac (16-84\%: -0.118 to +0.045), with 42\% of systematic Monte Carlo simulations showing a shortfall.
\begin{figure}\centering\includegraphics[width=\columnwidth]{result.png}\caption{The reionization ionizing-photon-budget shortfall for star-forming galaxies is not robust to systematics at z\~{}5 (lowxi systematic corner)}\end{figure}
\section{Caveats}
Caveats: Our results are subject to several limitations, primarily stemming from the reliance on automated, single-selection, uncalibrated measurements. The accuracy of our calculations depends heavily on the assumptions and calibrations used in previous studies, which may not fully capture the complexity of reionization processes. Additionally, our analysis does not account for potential variations in ξion or clumping factor C across different galaxy populations, which could introduce further uncertainties. Finally, the use of published literature values means that we cannot assess the impact of new observational data on our conclusions.
\section*{References}
\footnotesize
[Muoz2024] Muñoz et al. (2024). Reionization after JWST: a photon budget crisis?. bibcode 2024MNRAS.535L..37M \par [Park2022] Park et al. (2022). Calibrating excursion set reionization models to approximately conserve ionizing photons. bibcode 2022MNRAS.517..192P \par [Duncan2015] Duncan et al. (2015). Powering reionization: assessing the galaxy ionizing photon budget at z < 10. bibcode 2015MNRAS.451.2030D \par [Davies2021] Davies et al. (2021). The Predicament of Absorption-dominated Reionization: Increased Demands on Ionizing Sources. bibcode 2021ApJ...918L..35D \par [Madau2017] Madau (2017). Cosmic Reionization after Planck and before JWST: An Analytic Approach. bibcode 2017ApJ...851...50M

\end{document}
